% Grenzwerte und Verhalten im Unendlichen · Produkte aus Polynom und e-Funktion · Prüfungs-Fokus Abitur GK – bank/grenzwerte-und-verhalten-im-unendlichen/e2.jsonl (zusammenbau v0.6, Prüfungs-Fokus)
\einheitenkopf[e2]{Einheit 2 · Produkte aus Polynom und e-Funktion}
\zweigzeile{Abi GK ×7}
\setcounter{aufgabe}{0}

% Anlauf (ohne Prüfkennung)
\begin{aufgabe}{Produkte aus Polynom und e-Funktion – Anlauf}
\begin{teile}
\teil Welche Seite? Für welche Richtung geht bei $f(x) = (x - 6) \cdot e^{x}$ der e-Faktor gegen null? \kreuz{für $x \to +\infty$} \\ \kreuz{für $x \to -\infty$}
\teil $f(x) = e^{3x}$. Verhalten für $x \to +\infty$ und $x \to -\infty$? x → +∞: f(x) → \leerfeld ; x → −∞: f(x) → \leerfeld
\end{teile}
\end{aufgabe}

\begin{aufgabe}{Produkte aus Polynom und e-Funktion – Prüfungsaufgaben}
\begin{teile}
\teil $f(x) = (4 - x) \cdot e^{x}$. Gib das Verhalten für $x \to -\infty$ und $x \to +\infty$ an und sage, von welcher Seite sich der Graph der $x$-Achse nähert. \hfill (A25) \\ x → +∞: f(x) → \leerfeld ; x → −∞: f(x) → \leerfeld
\teil $f(x) = (x + 6) \cdot e^{x}$. Gib das Verhalten für $x \to -\infty$ und $x \to +\infty$ an und sage, von welcher Seite sich der Graph der $x$-Achse nähert. \hfill (A25) \\ x → +∞: f(x) → \leerfeld ; x → −∞: f(x) → \leerfeld
\teil $f(x) = (1 - 2x) \cdot e^{x}$. Gib das Verhalten für $x \to -\infty$ und $x \to +\infty$ an und sage, von welcher Seite sich der Graph der $x$-Achse nähert. \hfill (A25) \\ x → +∞: f(x) → \leerfeld ; x → −∞: f(x) → \leerfeld
\teil $f(x) = (x + 4) \cdot e^{-x}$. Grenzwert für $x \to +\infty$? Beschreibe den Verlauf des Graphen dort. \hfill (A22) \\ Grenzwert: \leerfeld
\teil $f(x) = (2x - 1) \cdot e^{-0{,}5x}$. Verhalten für $x \to +\infty$? Beschreibe den Verlauf des Graphen dort. \hfill (A18) \\ Grenzwert: \leerfeld
\teil $f(x) = (3x + 9) \cdot e^{-x}$. Grenzwert für $x \to +\infty$? Beschreibe den Verlauf des Graphen dort. \hfill (A22) \\ Grenzwert: \leerfeld
\teil $f(x) = 0{,}5 \cdot (x^2 - 9) \cdot e^{x}$. Verhalten für $x \to +\infty$ und $x \to -\infty$? \hfill (A23) \\ x → +∞: f(x) → \leerfeld ; x → −∞: f(x) → \leerfeld
\teil $f(x) = (-0{,}2x^2 + x) \cdot e^{-0{,}5x}$. Verhalten für $x \to +\infty$ und $x \to -\infty$? \hfill (A21) \\ x → +∞: f(x) → \leerfeld ; x → −∞: f(x) → \leerfeld
\teil $f(x) = 0{,}2x \cdot (4 - x) \cdot e^{x}$. Beschreibe den Verlauf des Graphen für $x \to -\infty$ und begründe am Term. \hfill (A20)
\teil Nach einem Regenguss ist der Pegel eines Bachs $w(t) = 6t \cdot e^{-0{,}5t} + 40$ ($t$ in h, $w$ in cm). Wie verhalten sich die Werte für $t \to +\infty$, und was bedeutet das? \hfill (A19)
\teil Nach einer Mahlzeit ist der Blutzucker $b(t) = 30t \cdot e^{-t} + 90$ ($t$ in h, $b$ in mg/dl). Wie verhalten sich die Werte für $t \to +\infty$, und was bedeutet das? \hfill (A19)
\teil Nach dem Einschalten einer Heizung ist die Temperatur eines Rohrs $T(t) = 12t \cdot e^{-0{,}2t} + 18$ ($t$ in min, $T$ in °C). Wie verhalten sich die Werte für $t \to +\infty$, und was bedeutet das? \hfill (A19)
\end{teile}
\end{aufgabe}

\begin{aufgabe}{Produkte aus Polynom und e-Funktion – Prüfungsaufgaben – weiter}
\begin{teile}
\teil $f(x) = (x - 5) \cdot e^{0{,}5x}$. Nullstelle und Verhalten für $x \to -\infty$ und $x \to +\infty$? \hfill (A25) \\ x = \leerfeld ; x → +∞: f(x) → \leerfeld ; x → −∞: f(x) → \leerfeld
\teil $f(x) = (6 - 2x) \cdot e^{x}$. Nullstelle und Verhalten für $x \to -\infty$ und $x \to +\infty$? \hfill (A25) \\ x = \leerfeld ; x → +∞: f(x) → \leerfeld ; x → −∞: f(x) → \leerfeld
\teil $f(x) = (x + 7) \cdot e^{-0{,}5x}$. Nullstelle und Verhalten für $x \to -\infty$ und $x \to +\infty$? \hfill (A18) \\ x = \leerfeld ; x → +∞: f(x) → \leerfeld ; x → −∞: f(x) → \leerfeld
\end{teile}
\end{aufgabe}

\medskip
%% Merkkasten Einheit 2: 8 Zeilen, Vorgabe höchstens fünf (3.1) – ungekürzt
\uebersichtskasten{e-Funktion: $e^x$ geht für x → −∞ gegen null und für x → +∞ gegen +∞; $e^{(−x)}$ umgekehrt; bei $e^{(kx)}$ entscheidet das Vorzeichen von k, auf welcher Seite die Werte gegen null gehen. \\ Die e-Funktion gewinnt: In einem Produkt p(x) · $e^x$ aus Polynom und e-Funktion bestimmt die e-Funktion das Verhalten. Auf der Seite, auf der $e^x$ gegen null geht, geht das Produkt gegen null – der Graph nähert sich der x-Achse (von oben, wenn p dort positiv ist, von unten, wenn negativ). Auf der anderen Seite geht das Produkt gegen +∞ oder −∞, je nach Vorzeichen von p(x) dort. \\ f(x) = (x − 3) · $e^x:$ für x → −∞ gilt f(x) → 0, der Graph nähert sich von unten der x-Achse (x − 3 < 0); für x → +∞ gilt f(x) → +∞. \\ g(x) = (5 − x) · $e^{(−x)}:$ für x → +∞ gilt g(x) → 0 (von unten, 5 − x < 0); für x → −∞ gilt g(x) → +∞ (beide Faktoren positiv und wachsend). \\ Konstanter Summand: bei p(x) · $e^{(−kx)}$ + c bleibt c als Grenzwert stehen, wenn das Produkt verschwindet. \\ Nullstellen: $e^x$ ist nie null – die Nullstellen von p(x) · $e^x$ sind genau die Nullstellen von p. \\ f(x) = (x − 3) · $e^x$ hat die einzige Nullstelle 3. \\ Parameter: bei (x + a) · $e^{(−x)}$ oder k · p(x) · $e^x$ hängt das Vorzeichen auf der wachsenden Seite vom Parameter ab – Fallunterscheidung nach dem Vorzeichen von k.}
